\documentclass[10pt,conference,compsocconf]{IEEEtran}
\setlength{\parindent}{0pt}

% ------------------------------------------------------
% Packages
\usepackage{amsmath, amssymb, bm}
\usepackage{graphicx}
\usepackage{booktabs}
\usepackage{hyperref}
\usepackage{hyperref}
\usepackage{graphicx}	% For figure environment

% ------------------------------------------------------
\title{\textbf{Predicting Cardiovascular Disease Using Machine Learning}}
\author{Yassine Zaki \and Rokas Bertasius \and Filippo Tognina}
\date{October 2025}

% ------------------------------------------------------
\begin{document}
\maketitle

% ------------------------------------------------------
\begin{abstract}
Cardiovascular diseases (CVDs) remain a leading global cause of death, underscoring the need for early detection. This project applies logistic regression to predict the risk of Myocardial Infarction or Coronary Heart Disease (MICHD) using Behavioral Risk Factor Surveillance System (BRFSS) data. After preparing the data the model was trained with gradient descent and L2 regularization, then evaluated using ROC metrics. Results show strong predictive accuracy, highlighting the potential of interpretable machine learning for large-scale health monitoring.
\end{abstract}

%-------------------------------------------------------
\section{Introduction}
Cardiovascular diseases (CVDs) are a leading cause of mortality in developed nations {ref}, accounting for millions of largely preventable deaths each year. Myocardial Infarction and Coronary Heart Disease (MICHD) are among the most common and often fatal without early detection. Identifying high-risk individuals early can significantly improve prevention and outcomes.
This project aims to develop a machine learning model to predict MICHD risk using data from the Behavioral Risk Factor Surveillance System (BRFSS). Each record represents an individual’s health profile, and the task is to predict whether that person is likely to develop MICHD. The project applies logistic regresion, focusing on model interpretability and practical use for population-level health monitoring.

% ------------------------------------------------------
\section{Model and Methods}
% ------------------------------------------------------
\subsection{Dataset and preprocessing}
The BRFSS dataset is a large-scale, phone-based survey collecting self-reported health data across domains such as health status, healthcare access, chronic conditions, demographics, lifestyle, and preventive behaviors.
Dataset comprises a heterogeneous mix of binary, categorical, and numerical attributes including administrative variables.
To render the data suitable for logistic regression:
\begin{itemize}
    \item Administrative variables, including phone numbers and call information, were dropped.
    \item Categorical variables were transformed using one-hot encoding.
    \item Missing values were handled by adding indicator columns that mark missing entries, allowing the model to treat missingness itself as informative. 
    \item Strong 10:1 class imbalance (low MICHD prevalence) was addressed by introducing class weights. Misclassifications of the minority class were penalized more heavily, reducing bias toward the majority class.   
\end{itemize}
The dataset was randomly split into a training set (80\%) and test set (20\%). 

\subsection{Model Overview}
Project employs \textbf{Logistic Regression}, a robust, classical binary classifier. The model estimates the probability that an individual develops MICHD given an input feature vector $\bm{x}$:

\begin{equation}
P(y = 1 \mid \bm{x}) = \frac{1}{1 + e^{-(\bm{w}^\top \bm{x} + b)}}
\end{equation}

where $\bm{w}$ denotes the learned weight vector, $\bm{x}$ the input feature vector, and $b$ the bias term.

Model hyperparameters, including learning rate, regularization strength, class weighting, were optimized via grid search. The optimal decision threshold was selected using best F1 score. Model weights were learned using Gradient Descent (GD) to minimize the negative log-likelihood loss. To improve generalization, L2 regularization with coefficient $\alpha = 0.1$ was applied, penalizing large weights and preventing overfitting. Resulting in the loss function defined in Eq.~\ref{eq:loss}. Training used early stopping with a patience of 20 epochs, halting optimization when validation loss stopped improving. 

\begin{equation}
\label{eq:loss}
\mathcal{L}(\bm{w}) = -\frac{1}{N} \sum_{i=1}^{N} 
\left[ y_i \log(\hat{y}_i) + (1 - y_i) \log(1 - \hat{y}_i) \right]
+ \alpha \lVert \bm{w} \rVert_2^2
\end{equation}

where $\hat{y}_i$ denotes the predicted probability for sample $i$.

% ------------------------------------------------------
\section{Results}


\textbf{##########to be bulshited tomorow}


This ROC curve, fig.1, illustrates the trade-off between the True Positive Rate (TPR) and the False Positive Rate (FPR) for different decision thresholds of our classifier. The curve shows how well the model distinguishes between positive and negative classes. The red point marks the threshold (0.77) that yielded the best F1-score (0.394), representing the optimal balance between precision and recall. A higher area under the curve (AUC) indicates better overall classification performance.

\begin{figure}
    \centering
    \includegraphics[width=1\linewidth]{roc_curve.png}
    \caption{ROC Curve}
    \label{fig:placeholder}
\end{figure}



This bar, Fig.2, chart presents the relative importance of the features used in our predictive model. Features with higher importance values have a stronger influence on the model’s predictions. Age-related variables such as \textbf{AGE80} and \textbf{AGEG5YR} are the most significant predictors, indicating that age plays a major role in determining the risk of cardiovascular diseases. Other influential factors include \textbf{CHILDREN, TOLDHI2} (doctor-diagnosed high blood pressure), and \textbf{RFCHOL} (cholesterol status), reflecting both demographic and health-related contributions to disease prediction.

\begin{figure}
    \centering
    \includegraphics[width=1\linewidth]{feature_importance.png}
    \caption{Features}
    \label{fig:placeholder}
\end{figure}

This plot, fig.3, shows the evolution of the training and validation losses over the training iterations. The training loss remains relatively stable, while the validation loss decreases steadily, indicating that the model continues to improve its performance on unseen data. The consistent gap between the two curves suggests that the model generalizes reasonably well, with no clear signs of overfitting. Overall, the downward trend in validation loss demonstrates that the optimization process successfully minimized the negative log-likelihood and improved predictive accuracy.



\begin{figure}
    \centering
    \includegraphics[width=1\linewidth]{loss_plot.png}
    \caption{Negative Log-Likelihood loss}
    \label{fig:placeholder}
\end{figure}





% ------------------------------------------------------
\section{Discussion}
[To include interpretation of model coefficients, ROC analysis, and feature importance insights.]

some bulshit what was tried and did not work, also add it to methods
% ------------------------------------------------------
\section{Summary}
[Summarize the predictive performance and clinical implications.]

% ------------------------------------------------------
\section{Limitations}
[Summarize the predictive performance and clinical implications.]

% ------------------------------------------------------
\end{document}

